\section{HCP base model and dynamic descriptor}
\label{sec:base_model_dynamic}

\subsection{Finite-strain crystal state and explicit local closure}

The homogeneous crystal follows the multiplicative split
\begin{equation}
 \boldsymbol F=\boldsymbol F_e\boldsymbol F_p,
 \qquad \det\boldsymbol F_p=1,
 \qquad
 \boldsymbol L_p=\dot{\boldsymbol F}_p\boldsymbol F_p^{-1}
 =\sum_{\alpha=1}^{18}\dot\gamma^\alpha
 \boldsymbol s^\alpha\otimes\boldsymbol m^\alpha .
 \label{eq:dyn_multiplicative_split}
\end{equation}
The HCP geometry contains three basal $\langle a\rangle$, three prismatic $\langle a\rangle$ and twelve pyramidal $\langle c+a\rangle$ systems. Twinning is inactive in the present baseline. The 62 independent local coordinates are
\begin{equation}
 \boldsymbol a=
 [\boldsymbol\xi_8,\boldsymbol\rho_m^{18},
  \boldsymbol\rho_d^{18},\boldsymbol\gamma^{18}]^{\mathsf T},
 \label{eq:dyn_local_state}
\end{equation}
where $\boldsymbol\xi$ is an eight-coordinate chart of the isochoric plastic distortion, $\rho_m^\alpha$ and $\rho_d^\alpha$ are mobile and dipole dislocation densities, and $\gamma^\alpha$ are signed slips. Imposing the chart before linearization removes the redundant plastic-volume coordinate.

For slip system $\alpha$, forest density, slip resistance and mean free path are
\begin{align}
 \rho_f^\alpha
 &=\max\!\left[\sum_{\beta=1}^{18}
 h_{\alpha\beta}(\rho_m^\beta+\rho_d^\beta),\rho_{\min}\right],
 \label{eq:ft_forest_density}\\
 g^\alpha
 &=\tau_0^\alpha+a_T\mu(T)b^\alpha\sqrt{\rho_f^\alpha},
 \label{eq:ft_slip_resistance}\\
 \Lambda^\alpha
 &=\left[d_g^{-1}+\frac{\sqrt{\rho_f^\alpha}}{K_\Lambda}\right]^{-1}.
 \label{eq:ft_mean_free_path}
\end{align}
The resolved Mandel stress $\tau^\alpha$ is coupled to the micromorphic slip $\zeta^\alpha$ through
\begin{equation}
 \tau_{\rm eff}^\alpha
 =\tau^\alpha+H_\chi(\zeta^\alpha-\gamma^\alpha).
 \label{eq:ft_effective_slip_stress}
\end{equation}
On an active smooth branch, the thermally activated kinetic law implemented in the base model is
\begin{align}
 \dot\gamma^\alpha
 &=\rho_m^\alpha b^\alpha v_0^\alpha
 \exp\!\left[-\frac{\Delta F^\alpha}{k_{\rm B}T}
 \left\langle1-\bar\tau_\alpha^{\,p^\alpha}\right\rangle^{q^\alpha}\right]
 \operatorname{sign}(\tau_{\rm eff}^\alpha),
 \label{eq:ft_slip_rate}\\
 \bar\tau_\alpha
 &=\min\!\left(
 \frac{\langle|\tau_{\rm eff}^\alpha|-g^\alpha\rangle}
 {\tau_{\rm cut}^\alpha},1\right),
 \qquad
 \dot\gamma^\alpha=0\quad\text{if }|\tau_{\rm eff}^\alpha|\le g^\alpha .
 \label{eq:ft_normalized_stress}
\end{align}
Here $\langle x\rangle=\max(x,0)$. The active-set condition in Eq.~\eqref{eq:ft_normalized_stress}, the density floor and the activation cap are recorded before differentiation; a perturbation probe that changes one of these branches is rejected.

The continuous-rate counterparts of the implemented semi-implicit density update can be written
\begin{align}
 \dot\rho_m^\alpha&=S^\alpha-(F^\alpha+G^\alpha)\rho_m^\alpha,
 &S^\alpha&=\frac{|\dot\gamma^\alpha|}{b^\alpha\Lambda^\alpha},
 \label{eq:ft_mobile_density}\\
 \dot\rho_d^\alpha&=F^\alpha\rho_m^\alpha-(G^\alpha+C^\alpha)\rho_d^\alpha,
 \label{eq:ft_dipole_density}
\end{align}
where
\begin{align}
 d_{\rm check}^\alpha&=c_d b^\alpha,
 &d_{\rm hat}^\alpha&=min\!\left[
 \max\!\left(\frac{3\mu_0b^\alpha}{16\pi|\tau_{\rm eff}^\alpha|},
 d_{\rm check}^\alpha\right),\Lambda^\alpha\right],\\
 F^\alpha&=\frac{2(d_{\rm hat}^\alpha-d_{\rm check}^\alpha)_+}{b^\alpha}
 |\dot\gamma^\alpha|,
 &G^\alpha&=\frac{2d_{\rm check}^\alpha}{b^\alpha}|\dot\gamma^\alpha|,\\
 C^\alpha&=\nu_c\exp[-Q_c/(k_{\rm B}T)].
 \label{eq:ft_density_rates}
\end{align}
The effective crystal-plastic power and its registered thermal/nonthermal
partition are
\begin{equation}
 \mathcal P_{\rm cp}=\sum_{\alpha=1}^{18}
 \tau_{\rm eff}^\alpha\dot\gamma^\alpha\ge0,
 \qquad r_T=\beta\mathcal P_{\rm cp},
 \qquad \dot u_s=(1-\beta)\mathcal P_{\rm cp},
 \qquad 0\le\beta\le1.
 \label{eq:ft_power_partition}
\end{equation}
Here $u_s$ is a passive nonthermal configurational ledger. It is not identified
with a dislocation free energy, generates no hardening force and is not an
active coordinate of the perturbation generator. The numerical base path
audits the Mandel/slip power equality and the mechanical--microforce
driving-power ledger.

Equations~\eqref{eq:dyn_multiplicative_split}--\eqref{eq:ft_power_partition}, anisotropic thermoelasticity and the local micromorphic force define
\begin{equation}
 (\boldsymbol P,r_T,\boldsymbol q,\dot{\boldsymbol a})
 =\boldsymbol{\mathcal R}(\boldsymbol F,T,\boldsymbol\zeta,\boldsymbol a).
 \label{eq:dyn_constitutive_map}
\end{equation}
The perturbation operator uses its continuous, branch-locked tangent
\begin{equation}
 \mathrm d\boldsymbol{\mathcal R}
 =\boldsymbol{\mathcal R}_{F}:\mathrm d\boldsymbol F
 +\boldsymbol{\mathcal R}_{T}\,\mathrm dT
 +\boldsymbol{\mathcal R}_{\zeta}\,\mathrm d\boldsymbol\zeta
 +\boldsymbol{\mathcal R}_{a}\,\mathrm d\boldsymbol a,
 \label{eq:ft_constitutive_tangent_contract}
\end{equation}
not an endpoint algorithmic tangent. Equation~\eqref{eq:ft_constitutive_tangent_contract} is the precise interface by which anisotropic elasticity, slip, dislocation evolution, heating and micromorphic coupling enter the descriptor.

\subsection{Recoverable micromorphic energy, entropy production and loading path}

The recoverable micromorphic contribution is
\begin{align}
 \psi_\chi
 &=\frac12H_\chi\|\boldsymbol\zeta-\boldsymbol\gamma\|^2
 +\frac12\operatorname{Grad}\boldsymbol\zeta:
 \boldsymbol H_\nabla:\operatorname{Grad}\boldsymbol\zeta,
 &\boldsymbol H_\nabla&=\boldsymbol H_{\rm slip}+\boldsymbol H_{\rm Nye},\\
 \boldsymbol H_{\rm Nye}&=\mu\ell_N^2
 \boldsymbol C_N^{\mathsf H}\boldsymbol C_N .
 \label{eq:dyn_gradient_energy}
\end{align}
Thus the energetic microstress and higher-order stress are
\begin{equation}
 \boldsymbol\pi=H_\chi(\boldsymbol\zeta-\boldsymbol\gamma),
 \qquad
 \boldsymbol\xi=\frac{\partial\psi_\chi}
 {\partial\operatorname{Grad}\boldsymbol\zeta},
 \qquad
 \frac{\partial\psi_\chi}{\partial\boldsymbol\gamma}=-\boldsymbol\pi .
 \label{eq:ft_energetic_microstresses}
\end{equation}
The last identity explains the effective driving stress
$\boldsymbol\tau_{\rm eff}=\boldsymbol\tau+\boldsymbol\pi$: after the
recoverable micromorphic-energy rate is removed from mechanical slip power,
the distributable remainder is $\mathcal P_{\rm cp}$ in
Eq.~\eqref{eq:ft_power_partition}. Under periodic or homogeneous
microtraction boundary conditions, the local split gives
\begin{equation}
 \mathcal D_{\rm mat}=\mathcal P_{\rm cp}-\dot u_s
 =r_T\ge0,
 \qquad
 \dot{\mathcal S}_{\rm prod}
 =\frac{r_T}{T}
 +\frac{\operatorname{Grad}T\cdot\boldsymbol\kappa
 \operatorname{Grad}T}{T^2}\ge0 .
 \label{eq:ft_entropy_production}
\end{equation}
The associated global first-law ledger contains kinetic energy, thermoelastic
internal energy, $\psi_\chi$ and $u_s$, with external mechanical power,
microtraction power and boundary heat flux. This closes the declared
power/heat bookkeeping, but it deliberately does not claim that $u_s$ is a
recoverable dislocation energy or that the density variables possess an
identified Helmholtz potential. A constitutively coupled dislocation-storage
law remains outside the present model.

The slip-gradient part of $\boldsymbol H_\nabla$ is full rank for the adopted
18-slip geometry; the Nye part is rank deficient and cannot alone control
every micromorphic direction. In the reference configuration,
\begin{align}
 \rho_0\ddot{\boldsymbol u}-\operatorname{Div}\boldsymbol P&=\boldsymbol0,
 \label{eq:dyn_balance_momentum}\\
 \rho_0c\dot T-\operatorname{Div}(\boldsymbol\kappa\operatorname{Grad}T)-r_T&=0,
 \label{eq:dyn_balance_heat}\\
 \boldsymbol q-\operatorname{Div}
 (\boldsymbol H_\nabla\operatorname{Grad}\boldsymbol\zeta)&=\boldsymbol0,
 \label{eq:dyn_balance_microforce}\\
 \dot{\boldsymbol a}-\boldsymbol f_a(\boldsymbol F,T,\boldsymbol\zeta,\boldsymbol a)&=\boldsymbol0.
 \label{eq:dyn_balance_state}
\end{align}
The homogeneous base path is simple shear,
\begin{equation}
 \boldsymbol F_0(t)=\boldsymbol I+\gamma_0(t)
 \boldsymbol e_1\otimes\boldsymbol e_2,
 \qquad \dot\gamma_0=10^4~\mathrm{s}^{-1},
 \qquad T_0(0)=293.15~\mathrm K,
 \label{eq:dyn_base_shear}
\end{equation}
integrated to $\gamma_0=1.4$ at $140~\mu\mathrm s$. It is a reproducible single-crystal research baseline, not a specimen-specific calibration.

\subsection{Fourier-linearized descriptor}

At a frozen base time, use
\begin{equation}
 \delta\boldsymbol y(\boldsymbol X,t)=\widehat{\boldsymbol y}
 \exp\!\left(s t+\mathrm ik\boldsymbol n\cdot\boldsymbol X\right),
 \qquad k>0,\quad\boldsymbol n\in\mathbb S^2 .
 \label{eq:dyn_fourier_ansatz}
\end{equation}
In the ordering $[\widehat{\boldsymbol u}_3,\widehat T,
\widehat{\boldsymbol\zeta}_{18},\widehat{\boldsymbol a}_{62}]$, linearization gives
\begin{equation}
 \underbrace{[\boldsymbol K(t;k,\boldsymbol n)+s\boldsymbol E(t)
 +s^2\boldsymbol M(t)]}_{\boldsymbol{\mathcal Q}(s;t,k,\boldsymbol n)}
 \widehat{\boldsymbol z}_{84}=\boldsymbol0.
 \label{eq:dynamic_descriptor}
\end{equation}
Only the three mechanical coordinates carry second-order inertia; temperature and the 62 local coordinates carry first-order storage; the 18 micromorphic coordinates are algebraic. The $\mathrm ik$ and $k^2$ blocks retain acoustic, thermomechanical, conduction and gradient couplings. Equation~\eqref{eq:dynamic_descriptor} is therefore a quadratic descriptor, not a conventional eigenproblem.
